%Version 3.1 December 2024
% See section 11 of the User Manual for version history
%
%%%%%%%%%%%%%%%%%%%%%%%%%%%%%%%%%%%%%%%%%%%%%%%%%%%%%%%%%%%%%%%%%%%%%%
%%                                                                 %%
%% Please do not use \input{...} to include other tex files.       %%
%% Submit your LaTeX manuscript as one .tex document.              %%
%%                                                                 %%
%% All additional figures and files should be attached             %%
%% separately and not embedded in the \TeX\ document itself.       %%
%%                                                                 %%
%%%%%%%%%%%%%%%%%%%%%%%%%%%%%%%%%%%%%%%%%%%%%%%%%%%%%%%%%%%%%%%%%%%%%

\documentclass[pdflatex,sn-mathphys-num]{sn-jnl}% Math/Physical Sciences Numbered style, single-column
                                                 % per editorial technical-check request

%%%% Standard Packages
\usepackage{graphicx}%
\usepackage{multirow}%
\usepackage{amsmath,amssymb,amsfonts}%
\usepackage{amsthm}%
\usepackage{mathrsfs}%
\usepackage{xcolor}%
\usepackage{textcomp}%
\usepackage{manyfoot}%
\usepackage{booktabs}%
\usepackage{tabularx}%
\usepackage{algorithm}%
\usepackage{algorithmicx}%
\usepackage{algpseudocode}%
\usepackage{listings}%
%%%%

\theoremstyle{thmstyleone}%
\newtheorem{theorem}{Theorem}%
\newtheorem{proposition}[theorem]{Proposition}%

\theoremstyle{thmstyletwo}%
\newtheorem{example}{Example}%
\newtheorem{remark}{Remark}%

\theoremstyle{thmstylethree}%
\newtheorem{definition}{Definition}%

\raggedbottom
\renewcommand{\arraystretch}{0.82}
\setlength{\textfloatsep}{4pt plus 1pt minus 1pt}
\setlength{\floatsep}{3pt plus 1pt minus 1pt}
\setlength{\intextsep}{3pt plus 1pt minus 1pt}
\setlength{\abovecaptionskip}{2pt}
\setlength{\belowcaptionskip}{1pt}
\AtBeginDocument{\setlength{\bibsep}{0pt}}
\renewcommand{\bibfont}{\footnotesize}

\begin{document}

\title[ARMOR-FL]{ARMOR-FL: Anytime-Valid Robust Martingale-Based Online Reweighting for Federated Intrusion Detection}

\author*[1]{\fnm{Quang-Vinh} \sur{Dang}}\email{vinh.dq4@buv.edu.vn}

\affil*[1]{\orgname{British University Vietnam}, \orgaddress{\city{Hung Yen}, \country{Vietnam}}}

\abstract{Federated intrusion detection lets network operators train a shared classifier without pooling raw traffic, but it opens the aggregator to poisoned client updates. Existing Byzantine-robust defenses compare each round's updates against the population and flag outliers, yet an anytime-valid comparison of this kind is mathematically guaranteed to eventually reject any client whose behavior sits persistently off-center -- so it cannot tell a merely non-IID honest client from an attacker without further information. We diagnose this failure concretely: a population-only e-process excluded up to 62.5\% of honest clients under IID data in our experiments, the opposite of what intuition suggests. We introduce ARMOR-FL, which adds a self-referential shift test -- has this client's own deviation changed, not just how large is it -- alongside a scale-robust population statistic and a threshold-calibrated trust weight, to recover the sequential-testing guarantee without this failure mode. Across CICIDS2017, CICIDS2018, and CICIoT2023 under label-flipping, Gaussian-noise, and adaptive (ALIE) attacks, ARMOR-FL removes the false-exclusion pathology entirely and outperforms undefended averaging, though a distance-based Krum baseline remains stronger in aggregate; we report this honestly, including the specific conditions -- non-IID data under lighter poisoning (a moderate 20\% malicious client fraction) -- where ARMOR-FL's mechanism gives it a real edge over Krum. We release the full implementation, configuration grids, and reassembly-only dataset archives.}

\keywords{Federated learning, Byzantine robustness, intrusion detection, anytime-valid inference, sequential testing, non-IID data}

\maketitle

\section{Introduction}\label{sec1}

A network intrusion detection system trained on one organization's traffic sees only that organization's attacks, but the traffic logs that would enrich it -- addresses, timing, payload sizes -- are exactly what no organization wants to hand over. Federated learning offers a way around this: each site trains locally and shares only model updates, and a central aggregator combines them into a global classifier none of the participants could have trained alone \citep{mcmahan2017fedavg}. This convenience creates an attack surface: an aggregator that averages every update by default cannot tell a legitimate one from one crafted to corrupt the global model, and unlike a centrally-trained system, an attacker here needs only federation membership, not data access. This is the Byzantine-robustness problem, and it is especially consequential for intrusion detection, where an attacker has a strong incentive to degrade the very system meant to catch it.

The dominant family of defenses treats each round's updates as a small sample and flags whichever deviates enough from the rest. Krum keeps only the most central update by nearest-neighbor distance \citep{blanchard2017krum}; FoolsGold tracks cosine similarity over time to catch colluding attackers \citep{fung2020foolsgold}; trimmed mean and coordinate-wise median discard per-coordinate extremes. All of these compare a client against the current round's \emph{population} -- and federated intrusion detection is a canonical non-IID setting, where different sites see different topologies, device populations, and background traffic, so honest updates are rarely interchangeable to begin with.

Our starting point was a puzzle that surfaced while building a defense for exactly this problem. Anytime-valid sequential testing -- e-processes and testing-by-betting \citep{waudbysmith2020betting} -- lets a defender accumulate evidence about a client across rounds rather than deciding in isolation each round, with a formal guarantee (Ville's inequality) on the false-exclusion rate at any stopping time. We built this, called it ARMOR-FL, and found that against synthetic honest clients with no attacker present at all, it excluded them anyway, sometimes catastrophically.

The reason is structural, not an implementation bug. An anytime-valid test of ``is this client's distance from the population median elevated'' is, by the very guarantee that makes it anytime-valid, mathematically certain to eventually reject any client whose true distance sits persistently above that median, whether the elevation reflects an attack or simply a client whose local data looks different from everyone else's. A defense built this way is guaranteed to fail, given enough time, for any client that never converges toward the population's center -- an ordinary heterogeneous participant as much as a persistently biased attacker. We confirmed this is not an edge case: in a controlled test, five of ten honest clients configured to be moderately more variable than the rest were excluded within twenty rounds by a population-only e-process. On real traffic it was worse: under CICIDS2017 with a purely IID partition and 20\% malicious clients -- the setting we expected to be \emph{easiest}, since IID data has no legitimate heterogeneity to confuse with attack -- the uncorrected defense excluded 62.5\% of honest clients. Homogeneity made the statistic \emph{more} fragile: with less genuine spread, the population's median-absolute-deviation shrinks, and dividing ordinary training noise by a near-zero scale manufactures large deviation scores from nothing unusual at all.

This paper reports the resulting design, ARMOR-FL, together with the diagnostic process, since we think the process is as informative as the fix. ARMOR-FL keeps the population e-process but adds a self-referential test: has this client's own deviation \emph{changed} relative to its own recent history, rather than how large it currently is. A client that has always looked off-center produces no shift signal and is protected; a client whose deviation appears mid-run does produce one, and only the combination of population-level and self-referential evidence triggers hard exclusion. We pair this with two further corrections: a population scale estimate that cannot collapse below a fraction of its own recent history, and a trust-weighting rule recalibrated against the same threshold that governs exclusion, so ordinary deviation is not punished disproportionately.

We are explicit that this is not a paper reporting a method that wins on every axis. Across CICIDS2017, CICIDS2018, and CICIoT2023, under label-flipping, Gaussian-noise, and adaptive (ALIE) attacks, ARMOR-FL removes the false-exclusion pathology it was built to fix and clearly outperforms undefended averaging, but Krum remains the stronger baseline in aggregate on all three datasets. What ARMOR-FL delivers is a specific regime -- non-IID populations under lighter label-flipping poisoning -- where its mechanism gives a genuine edge, together with a formal anytime-valid guarantee that Krum and FoolsGold do not provide, and an honest account of where it still falls short.

Our contributions are threefold. First, we identify and formally characterize a failure mode shared by population-only Byzantine-robust defenses under non-IID data: any anytime-valid test of population deviation alone is guaranteed to eventually exclude persistently off-center honest clients, worst, not best, under low heterogeneity. Second, we introduce ARMOR-FL, combining a self-referential shift test, a scale-floored population statistic, and a threshold-calibrated trust weight, validated on synthetic reproductions and three real datasets. Third, we provide, to our knowledge, the first reported robustness evaluation of a Byzantine-robust federated defense on CICIoT2023 under an explicit poisoning-attack grid, with a full open implementation, configuration files, and reassembly-only dataset archives.

\section{Related Work}\label{sec2}

\subsection{Byzantine-robust aggregation in federated learning}

The starting point for distance-based Byzantine robustness is Krum \citep{blanchard2017krum}, which scores each client by summed distance to its nearest neighbors and keeps only the most central update, discarding the rest. Multi-Krum generalizes this to a small committee, and Bulyan \citep{mhamdi2018bulyan} layers a coordinate-wise trimmed-mean filter on top to close a residual vulnerability to colluding attackers. Coordinate-wise trimmed mean and coordinate median trade Krum's winner-take-all selection for a per-coordinate robust average. FoolsGold instead exploits Sybil clients whose updates correlate over time and downweights them \citep{fung2020foolsgold}. Divide-and-Conquer (DnC) \citep{shejwalkar2021manipulating} counters adaptive ALIE-family attacks via spectral projection rather than distance, a natural comparison point for our own ALIE results (Section~\ref{sec:results}) that we do not evaluate directly, since our contribution is the sequential-testing layer rather than a new per-round filter. Karimireddy, He, and Jaggi \citep{karimireddy2021learning} diagnose a related tension -- heterogeneity inflates the same cross-sectional variance robust-aggregation rules use to detect attackers, degrading their convergence guarantee -- and propose worker momentum as a fix; their diagnosis targets the aggregation rule's convergence under heterogeneity in general, whereas ours targets a different failure of anytime-valid \emph{testing} rules layered on top of any such rule, and the two are complementary. All of these methods decide using only the current round's cross-sectional spread, which is exactly what makes them vulnerable to conflating honest heterogeneity with attack: a consistently-but-honestly different client looks, from a single round, identical to a consistently malicious one.

More recent work targets federated intrusion detection specifically: loss-scoring with Manhattan-similarity clustering \citep{yang2023dependable}; a systematic label-flipping study finding over ten points of accuracy loss in a well-tuned FL-IDS pipeline \citep{lavaur2025labelflip}, consistent with our own undefended-averaging results; graph-based clustering with Krum-style selection on IoT data \citep{rezaei2025grafids,rezaei2025rnn}; a Weibull-fit deviation model adjacent to our robust-scale statistic but without a sequential guarantee \citep{sameera2025weidetect}; gradient projection with isolation-forest scoring \citep{zukaib2026d3flids}; density-based collusion detection \citep{latif2026feddbc}; and an explainable defense at up to 30\% malicious participation, our own upper bound \citep{chowdhury2026flexids}. The broader \emph{Cluster Computing} literature is also active on adjacent problems \citep{akter2026fair,buyuktanir2025flidssurvey,mohawesh2026llmhfl,bella2026cpsfedids,godavarthi2026fedqlearning,houichi2026trustaware,ram2026blocshield}, but none combines Byzantine-robust aggregation with an anytime-valid statistical guarantee.

Two papers are closest in spirit to ARMOR-FL. Our own prior work, FORTRESS-FL \citep{dang2026fortressfl}, is a Byzantine-robust, privacy-preserving federated orchestration framework with no anytime-valid component and no mechanism to decouple drift from attack -- precisely the gap ARMOR-FL closes. RPCFL \citep{chen2026rpcfl}, published in this journal, is a Byzantine-robust, privacy-preserving clustered federated learning framework built on secure multiparty computation; it is, to our knowledge, the only other Byzantine-robust federated learning paper in \emph{Cluster Computing}. Its robustness guarantee is cryptographic rather than statistical and does not provide anytime-valid false-exclusion control, so the two approaches are complementary; a system combining SMPC-based privacy with anytime-valid testing is a natural future direction.

\subsection{Qualitative comparison of defense properties}\label{sec:qualcomp}

Table~\ref{tab:qualcomp} summarizes the eight defenses surveyed above (plus ARMOR-FL) along four properties that this paper's contribution turns on: whether the defense provides a formal, cumulative statistical guarantee rather than a per-round heuristic; whether it has any mechanism to decouple stable, honest heterogeneity from an evolving attack signal; its dominant per-round computational cost in the notation of Section~\ref{sec:complexity}; and whether it specifically targets coordinated (Sybil-style) attackers as opposed to independent ones.

\begin{table*}[htbp]
\centering
\caption{Qualitative comparison of Byzantine-robust federated defenses}\label{tab:qualcomp}
\footnotesize
\begin{tabularx}{\linewidth}{@{}XXXp{1.9cm}>{\centering\arraybackslash}p{1.6cm}@{}}
\toprule
Defense & Statistical guarantee & Drift/attack decoupling & Dominant cost & Sybil-targeted \\
\midrule
FedAvg (undefended) & None & N/A & $O(Kd)$ & No \\
Krum \citep{blanchard2017krum} & None (empirical) & None & $O(K^2d)$ & No \\
Bulyan \citep{mhamdi2018bulyan} & None (empirical) & None & $O(K^2d)$ & No \\
Trimmed mean / coord.\ median & None (empirical) & None & $O(Kd\log K)$ & No \\
FoolsGold \citep{fung2020foolsgold} & None (empirical) & None & $O(K^2)$ & Yes \\
DnC \citep{shejwalkar2021manipulating} & None (empirical) & None & $O(Kd)$ (spectral) & No \\
Karimireddy et al.\ \citep{karimireddy2021learning} & None (empirical) & Worker-momentum smoothing & $O(Kd)$ & No \\
FORTRESS-FL \citep{dang2026fortressfl} & None (empirical) & None & Not reported & No \\
RPCFL \citep{chen2026rpcfl} & None (cryptographic) & None & SMPC-dependent & No \\
ARMOR-FL (ours) & Anytime-valid ($\alpha$-controlled) & Self-shift e-process & $O(Kd)$ & No \\
\bottomrule
\end{tabularx}
\end{table*}

Three observations follow. First, ARMOR-FL is, to our knowledge, the only defense here offering a formal cumulative false-exclusion guarantee, but the guarantee is only as good as the population statistic it is built on, and Section~\ref{sec:results} shows it does not translate into an aggregate accuracy advantage over Krum. Second, we distinguish two senses of ``decoupling heterogeneity from attack'': Karimireddy et al.\ \citep{karimireddy2021learning} smooth the \emph{aggregation rule itself} via worker momentum so its convergence guarantee survives heterogeneity, whereas ARMOR-FL addresses a different failure specific to \emph{sequential exclusion decisions} -- whether a test accumulating evidence across rounds can tell persistent-but-stable deviation from attack-coincident deviation; no other defense here provides this second form, and the two mechanisms are complementary. Third, RPCFL's guarantee is cryptographic rather than statistical, complementary to rather than a substitute for ARMOR-FL's approach.

\subsection{Drift, heterogeneity, and the limits of population comparison}

The tension between non-IID heterogeneity and outlier-based robustness has been noted before, but rarely diagnosed at the level of a specific, quantified failure mode. Our own earlier work, ST-FedXIDS \citep{dang2026stfedxids}, addresses drift adaptation via graph learning for high-density WiFi environments but treats drift as something to adapt to rather than formally decouple from an attack test; we cite it for the drift-adaptation half of the present contribution only. Uncertainty-measure approaches to intrusion detection \citep{dang2026uncertainty} motivate the statistical-testing framing we build on, without addressing the federated non-IID setting directly.

The broader non-IID federated learning literature responds architecturally -- personalization layers, similarity-based clustering, client-side regularization, or worker-momentum smoothing of the aggregation rule itself \citep{karimireddy2021learning} -- rather than by changing the statistical test a Byzantine defense uses to flag a client. None of these touches the sequential-exclusion mechanism specifically: to our knowledge, no prior published defense identifies that an anytime-valid population test is not merely inaccurate under heterogeneity but mathematically \emph{guaranteed} to eventually misclassify a persistently off-center honest client, with the failure growing worse, not better, as the population becomes more homogeneous -- against the intuitive expectation that IID data is the easy case.

\subsection{Anytime-valid inference and its (limited) use in federated defenses}

E-processes and testing-by-betting \citep{waudbysmith2020betting} generalize sequential testing to a valid Type-I error guarantee at every stopping time via Ville's maximal inequality for nonnegative supermartingales. Ramdas, Gr\"unwald, Vovk, and Shafer's survey of safe anytime-valid inference \citep{ramdas2023savi} situates this within a broader family of time-uniform guarantees that includes Howard, Ramdas, McAuliffe, and Sekhon's confidence sequences \citep{howard2021timeuniform}; neither has, to our knowledge, been previously applied to Byzantine exclusion. This theory has seen substantial uptake for continuous-monitoring problems -- clinical-trial interim analysis, A/B testing, drift detection -- but limited adoption inside federated Byzantine defenses, all of which make an independent per-round decision with no cumulative false-exclusion guarantee. ARMOR-FL's fixed, theoretically grounded null mean, rather than online calibration from a client's own early behavior (which would let an already-compromised client normalize itself), is a design choice we believe is necessary: an online-calibrated null risks exactly the self-normalization failure it is meant to prevent.

\subsection{Dataset coverage}

CICIDS2017 \citep{sharafaldin2018cicids} and its successor CSE-CIC-IDS2018 remain the most widely used benchmarks in the Byzantine-robust federated learning literature surveyed above. CICIoT2023 \citep{neto2023ciciot2023}, a larger, more recent IoT-focused dataset, has seen substantial adoption for centralized intrusion detection \citep{wahab2025multiclass,wakili2025resilient} but, to our knowledge, no prior Byzantine-robust federated defense with an explicit poisoning-attack evaluation has reported results on it -- a first-mover contribution here, since ARMOR-FL's markedly weaker aggregate performance on CICIoT2023 than on the two CIC datasets (Table~\ref{tab:master}) is itself new information, discussed further in Section~\ref{sec:discussion}.

The present work builds directly on FedSE-1DSqueezeNet \citep{zhou2026fedse1dsqueezenet}, published in this journal, whose backbone architecture we reimplement without modification (Section~\ref{sec:setup}); our contribution is entirely at the aggregation layer, and we adopt its federated hyperparameters (Table~\ref{tab:hyperparams}) as closely as computational constraints allow. Recent centralized work on these same datasets \citep{xu2025fewshot,zhang2025transformer,palani2026comparative,xue2025haehrl,sireesha2026lightweight,zhao2024gan} does not address federated poisoning robustness but establishes that the classification task is solvable to a high standard when data can be pooled centrally, useful context for the accuracy ceilings in Section~\ref{sec:results}.

Our own prior work spans federated IoT malware detection \citep{dang2024iotmalwarefl}, conformal-prediction botnet detection \citep{dang2024conformalbotnet}, cross-dataset transfer \citep{dang2023transferids}, general IoT intrusion detection \citep{dang2022iotids,dang2022mlids}, and explainability-driven detection \citep{dang2021explainableids}; none combines Byzantine robustness with anytime-valid inference, the gap this work fills. A separate personalized federated poisoning defense \citep{thein2024pflids} similarly stops short of an anytime-valid guarantee.

\section{Method}\label{sec:method}

\subsection{Preliminaries: anytime-valid sequential testing}\label{sec:prelim}

Before describing ARMOR-FL's aggregation pipeline, we make precise the statistical machinery underlying it, since the failure mode motivating this paper (Section~\ref{sec1}) is a consequence of this machinery's own guarantee rather than an implementation defect.

\begin{definition}[e-process]\label{def:eprocess}
Let $(\mathcal{F}_t)_{t \ge 0}$ be a filtration representing the information available to the aggregator after round $t$. A process $(E_t)_{t \ge 0}$ adapted to $(\mathcal{F}_t)$ is an \emph{e-process} for a null hypothesis $H_0$ if $E_0 = 1$ and $(E_t)$ is a nonnegative supermartingale under $H_0$, i.e.\ $E_t \ge 0$ and $\mathbb{E}[E_t \mid \mathcal{F}_{t-1}] \le E_{t-1}$ for all $t \ge 1$.
\end{definition}

\begin{theorem}[Ville's inequality, anytime-valid Type-I control]\label{thm:ville}
Let $(E_t)$ be a nonnegative supermartingale with $E_0 = 1$ under $H_0$. Then for any $\alpha \in (0,1)$,
\[
\Pr_{H_0}\!\left[\exists\, t \ge 0 : E_t \ge 1/\alpha\right] \le \alpha .
\]
\end{theorem}

Theorem~\ref{thm:ville} is the classical result underlying testing-by-betting \citep{waudbysmith2020betting}; its proof (a consequence of the supermartingale maximal inequality) holds uniformly over \emph{every} stopping rule, including one that stops the first time $E_t$ crosses $1/\alpha$ -- precisely how the aggregator decides when to exclude a client. This is the appeal of the anytime-valid framework over fixed-sample testing: the defender need not commit in advance to how many rounds of evidence to collect, and the false-exclusion probability is controlled at level $\alpha$ no matter when that decision is made.

ARMOR-FL instantiates $(E_t)$ via the testing-by-betting construction: given a per-round statistic $z_k^{(t)}$ for client $k$ with null mean $\mathbb{E}[z_k^{(t)} \mid \mathcal{F}_{t-1}] = \mu_0$ under $H_0$ (honest behavior), the per-round multiplier is $e_t = 1 + \lambda(z_k^{(t)} - \mu_0)$ for a fixed betting fraction $\lambda \in (0,1)$, and $E_t = \prod_{s \le t} e_s$. Because $\mathbb{E}[e_t \mid \mathcal{F}_{t-1}] = 1$ under $H_0$, $(E_t)$ is exactly a nonnegative martingale with $E_0=1$, so Theorem~\ref{thm:ville} applies directly. ARMOR-FL fixes $\mu_0 = 1$ for the population statistic (Section~\ref{sec:method}.2) rather than calibrating it online from a client's own early behavior, which would let an already-compromised client's early rounds define its own null, defeating the guarantee entirely.

\begin{proposition}[Guaranteed eventual exclusion under persistent elevation]\label{prop:eventual-exclusion}
Suppose a client's true per-round statistic satisfies $\mathbb{E}[z_k^{(t)} \mid \mathcal{F}_{t-1}] = \mu > 1$ for every round $t$ (a client whose deviation from the population center is persistently, not merely transiently, elevated), and suppose in addition that $\mathrm{Var}(z_k^{(t)} \mid \mathcal{F}_{t-1}) \le \sigma^2 < \infty$ uniformly in $t$ -- a mild regularity condition satisfied whenever the client's round-to-round training dynamics are stable rather than diverging, which we take as given throughout (both honest heterogeneous clients and the attacks in our experimental grid produce stable, non-diverging per-round distances). Then for a suitable fixed betting fraction $\lambda \in (0,1)$, $\log E_t \to +\infty$ almost surely as $t \to \infty$, so for any significance level $\alpha \in (0,1)$ there exists an almost-surely finite round $T$ at which $E_T \ge 1/\alpha$, i.e.\ the client is eventually excluded with probability one.
\end{proposition}

\begin{proof}[Proof sketch]
Let $f(\lambda) = \mathbb{E}[\log(1 + \lambda(z_k^{(t)} - 1))]$ at the persistent mean $\mu$. Then $f(0)=0$ and $f'(0)=\mu-1>0$, so some $\lambda^*\in(0,1)$ gives $f(\lambda^*)>0$. Writing $X_t=\log e_t$, boundedness of $1+\lambda^*(z_k^{(t)}-1)$ below by $1-\lambda^*>0$ bounds the derivative of $\log(1+\lambda^*(z-1))$, giving $\mathrm{Var}(X_t\mid\mathcal{F}_{t-1})$ uniformly bounded from the assumed uniform variance bound on $z_k^{(t)}$. Kolmogorov's strong law for martingale differences then gives $\frac{1}{t}\sum_{s\le t}X_s \to f(\lambda^*)>0$ almost surely, so $\log E_t \to +\infty$ almost surely and crosses the fixed threshold $\log(1/\alpha)$ in finite time.
\end{proof}

Proposition~\ref{prop:eventual-exclusion} restates a known property of sequential tests -- power against a fixed persistent alternative tends to one as observations grow, traceable to Wald's sequential probability ratio test \citep{wald1945sprt} -- but applied to a population-comparison statistic as a Byzantine-exclusion gate, this consistency is precisely the failure mode: nothing distinguishes an attacker from a merely heterogeneous honest client, since both produce $\mu>1$ persistently, and both are eventually excluded with probability one. A defense built purely on Theorem~\ref{thm:ville} therefore needs an additional signal distinguishing \emph{why} $\mu>1$, exactly the role of ARMOR-FL's self-shift test (Section~\ref{sec:method}.3).

\begin{remark}[Scale-estimator fragility under low heterogeneity]\label{rem:mad-fragility}
A related fragility affects the scale estimate $\mathrm{MAD}^{(t)}$ itself: it is a median over only the $|\mathcal{S}_t|$ active clients sampled that round (six of ten here), so its relative error does not vanish as the population becomes more homogeneous -- if anything the opposite, since as true dispersion shrinks toward the noise floor of ordinary training variability, a client whose absolute deviation does not shrink at the same rate produces an inflated $z_k^{(t)}$ through no fault of its own. Section~\ref{sec:method}.2 shows this severe enough to produce a $62.5\%$ benign false-exclusion rate under IID CICIDS2017 data before correction, the empirical grounding for the MAD-flooring fix.
\end{remark}

\subsection{Problem setting and backbone model}

We consider a federated learning system with $K$ clients, of which an unknown subset $\mathcal{M} \subset \{1,\dots,K\}$ may be malicious. In each round $t$, the aggregator selects a subset $\mathcal{S}_t$ (client\_fraction $=0.6$ of $K=10$, i.e.\ six per round), broadcasts $\theta^{(t)}$, and each selected client $k$ trains locally for up to $E$ epochs with early stopping on patience $P$, producing $\theta_k^{(t)}$. The aggregator combines $\{\theta_k^{(t)} : k \in \mathcal{S}_t\}$ into $\theta^{(t+1)}$, suppressing malicious contributions without discarding legitimate heterogeneity.

Our backbone classifier is SE-1DSqueezeNet, reimplemented layer-for-layer from FedSE-1DSqueezeNet \citep{zhou2026fedse1dsqueezenet}: a 1D convolution, a depthwise-dilated-squeeze-excite module, global average pooling, and a fully connected head, over length-$F$ single-channel flow features. We reuse this backbone without modification, since our contribution is entirely at the aggregation layer; full architectural detail is in \citep{zhou2026fedse1dsqueezenet} and our released implementation.

\subsection{Baseline aggregation rules}

We compare ARMOR-FL against three baselines. \textbf{FedAvg} \citep{mcmahan2017fedavg} computes a sample-size-weighted average with no robustness mechanism, the undefended floor. \textbf{Krum} \citep{blanchard2017krum} scores each client by summed squared distance to its $k{-}f{-}2$ nearest neighbors ($f$ the assumed Byzantine count) and selects the single lowest-scoring update as the round's entire contribution. \textbf{FoolsGold} \citep{fung2020foolsgold} tracks cosine similarity between clients' update histories and downweights suspiciously similar contributions, targeting coordinated Sybil-style attacks rather than independent ones.

\subsection{ARMOR-FL aggregation}

ARMOR-FL's per-round pipeline (Algorithm~\ref{alg:armor}) proceeds in five stages, three of which are corrected forms of mechanisms whose naive versions we found, during development, to actively harm robustness under realistic non-IID conditions; we describe each naive version and the failure that motivated its replacement, since we consider the failure modes as reusable a contribution as the fixes.

\textbf{Robust reference center.} Given the active (non-excluded) clients' vectors, ARMOR-FL computes a coordinate-wise trimmed mean $m^{(t)}$ (trim fraction $0.2$) as the round's reference point, and each client's distance to it, $d_k^{(t)} = \lVert \theta_k^{(t)} - m^{(t)} \rVert$.

\textbf{Scale-floored population deviation.} The naive population statistic is $z_k^{(t)} = d_k^{(t)} / \mathrm{MAD}^{(t)}$, where $\mathrm{MAD}^{(t)}$ is the median of $\{d_j^{(t)}\}_{j \in \mathcal{S}_t}$ -- a robust scale estimate by construction. The failure is that $\mathrm{MAD}^{(t)}$ is estimated from only the current round's small, randomly-selected subset (six of ten here), and when the honest population happens to look unusually similar -- genuinely IID data, or simply a subsample that agreed closely -- $\mathrm{MAD}^{(t)}$ collapses toward zero, and dividing ordinary training noise by a near-zero denominator manufactures large $z_k^{(t)}$ values from nothing unusual. We confirmed this on CICIDS2017 under a purely IID partition with 20\% malicious clients: the uncorrected statistic excluded 62.5\% of honest clients, the highest false-exclusion rate anywhere in our grid, in the setting with the least genuine heterogeneity to explain it. ARMOR-FL floors the scale estimate at a fraction (\texttt{mad\_floor\_fraction}$=0.5$) of its own recent, unfloored history so the floor cannot ratchet upward on its own account: $\mathrm{MAD}^{(t)}_{\mathrm{eff}} = \max\bigl(\mathrm{MAD}^{(t)},\ 0.5 \cdot \mathrm{median}(\mathrm{MAD}^{(t-w)},\dots,\mathrm{MAD}^{(t-1)})\bigr)$. A genuine attack's raw MAD passes through unaffected; only spurious collapse is corrected.

Each client's population e-process is updated with $z_k^{(t)}$ via the fixed-mean betting construction of \citep{waudbysmith2020betting}: $e_t = 1 + \lambda(z_k^{(t)} - 1)$ with $\lambda = 0.5$ and null mean fixed at one (never recalibrated online, since an already-compromised client would otherwise learn to appear normal), and $\log E_{\mathrm{pop},k}^{(t)} = \sum_{s\le t} \log e_s$ accumulates evidence under Ville's inequality.

\textbf{Self-referential drift and shift tests.} A parallel e-process compares each client's local validation loss to the median and MAD of its own loss history, intended to decouple ``drifting'' from ``attacking.'' We found this insufficient alone: a client persistently elevated in $z_k^{(t)}$ due to ordinary heterogeneity typically has a perfectly \emph{stable} own-loss history, so the drift test never flags it and the population e-process excludes it anyway -- in a controlled synthetic test, all five of ten honest clients configured with roughly double the parameter-update variance of the rest were excluded within twenty rounds, despite no attack.

ARMOR-FL adds a second test on the client's own \emph{population-deviation history} rather than its loss: $\mathrm{shift}_k^{(t)} = z_k^{(t)} / \mathrm{median}(z_k^{(t-b)}, \dots, z_k^{(t-1)})$ over a baseline window $b=2$. A client off-center since the start produces a shift near one indefinitely and accumulates no shift evidence; a client whose deviation appears mid-run -- how every attack in our grid is configured (Section~\ref{sec:setup}) -- produces a detectable shift. Hard exclusion now requires \emph{both} population evidence past threshold \emph{and} either shift evidence or a gross-outlier deviation ($z_k^{(t)} \ge 4$): $\text{exclude}_k^{(t)} = \mathbb{1}\bigl[E_{\mathrm{pop},k}^{(t)} \ge 1/\alpha_{\mathrm{pop}}\bigr] \wedge \bigl(\mathbb{1}[\mathrm{shift\ evidence}] \vee \mathbb{1}[z_k^{(t)} \ge 4]\bigr)$. Re-running the synthetic test with this gate produced zero false exclusions across twenty-five rounds.

This gate trades one failure mode for a narrower one: a client Byzantine from round one that remains \emph{constant} thereafter, below the gross-outlier threshold, produces no shift signal either, since nothing changes relative to its own already-compromised baseline. This is a genuine information-theoretic limit, not a bug -- a constant, moderate-magnitude attacker and a heterogeneous-but-honest client are indistinguishable from behavior alone, and resolving it would require external information the aggregator does not have (Section~\ref{sec:limitations}).

\textbf{Threshold-calibrated trust weighting.} The naive rule, $w_k^{(t)} = n_k \cdot \mathrm{sigmoid}(-\beta \log E_{\mathrm{pop},k}^{(t)})$ with $\beta=4$, applies the sigmoid directly to the unbounded log-evidence. This disproportionately punishes mild deviation: at $\log E_{\mathrm{pop},k}^{(t)} = 1.0$, two orders of magnitude below the exclusion threshold ($-\log(0.05) \approx 3.0$), it already reduces weight to about 2\% of a client's share. On real CICIDS2017 data under label-flipping, malicious clients -- whose updates often looked statistically \emph{central} rather than extreme, a known blind spot of distance-based defenses -- retained near-full weight while honest clients with only ordinary deviation were crushed to a few percent, and ARMOR-FL performed worse than plain FedAvg as a direct result. ARMOR-FL instead rescales the sigmoid's input by the same threshold that gates hard exclusion, $\tau = -\log(\alpha_{\mathrm{pop}})$: $w_k^{(t)} = n_k \cdot \mathrm{sigmoid}\bigl(-\beta (\log E_{\mathrm{pop},k}^{(t)} / \tau - 1)\bigr)$, giving trust $\approx 0.98$ at the population median and exactly $0.5$ at the exclusion boundary -- the two reference points are fixed by the same $\alpha_{\mathrm{pop}}$ that governs hard exclusion, rather than an independently chosen constant on an unbounded raw statistic, and only evidence well past the boundary drops weight meaningfully.

\textbf{Probation and reinstatement.} An excluded client re-enters after a fixed probation window (three rounds) at reduced trust, with both e-processes reset -- a reset that gives a reinstated attacker a fresh accumulation window before re-exclusion, bounded by how quickly the population e-process can re-accumulate evidence (Section~\ref{sec:limitations}).

\begin{algorithm}
\caption{ARMOR-FL aggregation, one round}\label{alg:armor}
\footnotesize
\begin{algorithmic}[1]
\State \textbf{Input:} active client updates $\{\theta_k^{(t)}, n_k, \ell_k^{(t)}\}_{k \in \mathcal{S}_t}$, prior state
\State $m^{(t)} \gets$ coordinate-wise trimmed mean of $\{\theta_k^{(t)}\}$
\State $d_k^{(t)} \gets \lVert \theta_k^{(t)} - m^{(t)} \rVert$ for each $k$
\State $\mathrm{MAD}^{(t)}_{\mathrm{eff}} \gets \max\bigl(\mathrm{median}(d^{(t)}),\ 0.5 \cdot \mathrm{median}(\text{raw MAD history})\bigr)$
\For{each active client $k$}
  \State $z_k^{(t)} \gets d_k^{(t)} / \mathrm{MAD}^{(t)}_{\mathrm{eff}}$
  \State update population e-process with $z_k^{(t)}$
  \State update drift e-process with own-loss deviation
  \State update shift e-process with $z_k^{(t)} / \mathrm{median}(\text{own } z \text{ history})$
  \State $w_k^{(t)} \gets n_k \cdot \mathrm{sigmoid}\bigl(-\beta(\log E_{\mathrm{pop},k}^{(t)}/\tau - 1)\bigr)$
  \If{population evidence $\ge$ threshold \textbf{and} (shift evidence \textbf{or} $z_k^{(t)} \ge 4$)}
    \State exclude $k$; begin probation; $w_k^{(t)} \gets 0$
  \EndIf
\EndFor
\State $\theta^{(t+1)} \gets$ weighted average of $\{\theta_k^{(t)}\}$ with weights $\{w_k^{(t)}\}$
\State \textbf{Output:} $\theta^{(t+1)}$
\end{algorithmic}
\end{algorithm}

\subsection{Computational complexity}\label{sec:complexity}

Let $K$ be the number of active clients and $d$ the parameter dimension. FedAvg computes a single weighted mean, $O(Kd)$. ARMOR-FL's trimmed-mean reference center is $O(Kd \log K)$ with a per-coordinate sort, dominated in practice by $O(Kd)$; distance computation and e-process updates add $O(Kd)$ plus $O(K)$ bookkeeping, so ARMOR-FL's aggregation overhead over FedAvg is the same asymptotic order, $O(Kd)$. Krum requires the full pairwise distance matrix, $O(K^2d)$, the asymptotically most expensive of the four rules, though modest at our client counts ($K=10$, six active per round). FoolsGold additionally maintains a pairwise cosine-similarity history, $O(K^2)$ per round.

This ordering is reflected in wall-clock measurements (Table~\ref{tab:overhead}, Section~\ref{sec:results}.6): Krum is, in practice, the slowest of the four aggregators, not ARMOR-FL -- its aggregation-layer cost is small relative to local training time, which dominates total wall-clock time for all four rules. This is a relevant counterpoint to Section~\ref{sec:results}.1's headline finding that Krum wins on accuracy: the advantage does not come for free, and a latency-sensitive deployment may still prefer ARMOR-FL's cost profile even where its accuracy trails.

\section{Experimental Setup}\label{sec:setup}

\subsection{Datasets}

We evaluate on three widely used network intrusion detection datasets. \textbf{CICIDS2017} \citep{sharafaldin2018cicids} contains five days of labeled traffic covering benign activity and several attack types (DoS, DDoS, brute force, infiltration, web attacks, botnet) on a realistic testbed. \textbf{CICIDS2018} extends the same methodology across a larger ten-day capture. \textbf{CICIoT2023} \citep{neto2023ciciot2023} is a more recent IoT-focused dataset covering 33 attack types across seven categories from a 105-device smart-home testbed; we are not aware of any prior Byzantine-robust federated defense reporting a poisoning-attack evaluation on it. All datasets are partitioned across ten simulated clients using either an IID split or a Dirichlet-distributed non-IID split with concentration $\alpha \in \{0.5, \mathrm{IID}\}$ (smaller $\alpha$ is more skewed). Reflecting the practical compute budget, we sample 10\% of CICIDS2017, 2\% of CICIDS2018, and 1\% of CICIoT2023, calibrated so CICIDS2017's run-time sets the sampling fraction for the other two; resulting training sets remain in the hundreds of thousands of rows.

\subsection{Attacks}\label{sec:attackimpl}

We evaluate three attack types of increasing subtlety, beginning at round three (\texttt{attack\_start\_round}$=3$) after a two-round clean baseline, so every client has an honest baseline period before any attack, which is what makes ARMOR-FL's shift signal meaningful. We test malicious fractions of 0.2 and 0.4 (two and four of ten clients). \textbf{Label-flipping} applies a deterministic class shift, $y \mapsto (y + 1) \bmod C$ for $C$ classes, before training; no perturbation is applied to the resulting update, which is why it remains statistically close to an honest one (Section~\ref{sec:method}) -- a blind spot shared by every distance-based defense we test. \textbf{Gaussian-noise} replaces the malicious update entirely with Gaussian noise, norm-matched to what an honest update would have had: $\Delta_{\mathrm{malicious}} = \sigma \cdot \|\Delta_{\mathrm{honest}}\| \cdot \eta / \|\eta\|$ for $\eta \sim \mathcal{N}(0, I)$, $\sigma=1.0$; norm-matching evades a trivial magnitude check and, we believe (Section~\ref{sec:case2}), is also why the attack proved catastrophic on CICIoT2023 -- a norm-matched random-direction replacement discards all honest gradient signal while preserving scale, more destructive on a harder task where more of the update's norm carries genuine signal. \textbf{ALIE} (a little is enough) \citep{baruch2019alie} computes the coordinate-wise mean and standard deviation of the round's honest updates and returns $\Delta_{\mathrm{malicious}} = \bar{\Delta}_{\mathrm{honest}} - z_{\max} \cdot \sigma_{\mathrm{honest}}$ with $z_{\max}=1.5$, an omniscient construction requiring visibility into other clients' updates the same round, standard in the ALIE literature, designed to stay within the honest population's statistical spread and evade distance-based detection.

\subsection{Hyperparameters and deviations from the original protocol}\label{sec:hyperparams}

Table~\ref{tab:hyperparams} lists the hyperparameters used. We follow FedSE-1DSqueezeNet's protocol \citep{zhou2026fedse1dsqueezenet} where feasible, with one disclosed deviation: batch size. The original protocol specifies 32; a controlled sweep on identical data found per-round wall-clock time dropping from 211.5s at batch size 32 to 25.9s at 256, an eightfold reduction, with no measurable accuracy or F1 change across batch sizes 32--512. We attribute this to the backbone being small and the GPU large enough that batch size 32 is latency- rather than compute-bound (GPU utilization measured 5--18\%). We use batch size 256 throughout the robustness grid on this basis, a compute-budget decision disclosed rather than presented silently as the original setting. Separate comparability runs reproducing the original paper's own sanity check (R=5, FedAvg-only, IID plus three non-IID levels) retain batch size 32, for fidelity to the original protocol.

\begin{table*}[htbp]
\centering
\caption{Federated learning hyperparameters, robustness grid}\label{tab:hyperparams}
\footnotesize
\begin{tabular}{@{}ll@{}}
\toprule
Parameter & Value \\
\midrule
Clients ($K$) & 10 \\
Client fraction per round & 0.6 \\
Communication rounds & 25 \\
Local epochs (max, with early stopping) & 10 \\
Local early-stopping patience & 3 \\
Learning rate & 0.01 \\
Batch size & 256 (see Section~\ref{sec:hyperparams}) \\
Dropout & 0.5 \\
Attack start round & 3 \\
ARMOR-FL burn-in rounds & 3 \\
ARMOR-FL $\alpha_{\mathrm{pop}}$, $\alpha_{\mathrm{drift}}$, $\alpha_{\mathrm{shift}}$ & 0.05 each \\
ARMOR-FL trim fraction & 0.2 \\
ARMOR-FL MAD floor fraction & 0.5 \\
ARMOR-FL gross-outlier threshold ($z$) & 4.0 \\
ARMOR-FL probation rounds & 3 \\
\bottomrule
\end{tabular}
\end{table*}

\subsection{Evaluation metrics}\label{sec:metrics}

We report two accuracy statistics per run: \textbf{final-round accuracy}, comparable across all four aggregators, and \textbf{mean accuracy over the final eight rounds}, a more stable statistic since individual runs oscillate considerably -- a property of the random six-of-ten subsampling, not unique to ARMOR-FL. We additionally report the \textbf{stability gap}, $|\,\text{final-round accuracy} - \text{mean last-8-round accuracy}\,|$, summarizing how misleading a single final-round number can be.

For ARMOR-FL specifically, we report three detection statistics against ground-truth malicious labels (not available to the aggregator itself): \textbf{detection precision} (fraction of ever-excluded clients that are truly malicious), \textbf{detection recall} (fraction of truly malicious clients excluded at least once), and \textbf{benign false-exclusion rate} (fraction of truly honest clients excluded at least once). All three are averaged within a run and then across the twelve attack-fraction-$\alpha$ combinations per dataset, matching Table~\ref{tab:master}'s aggregation.

\subsection{Ablation protocol}\label{sec:ablation-protocol}

To isolate the effect of the three corrections in Section~\ref{sec:method} from the aggregation mechanism's design in general, we retain the full robustness grid from an earlier, uncorrected ARMOR-FL implementation, run under identical hyperparameters, datasets, attacks, and combinations. We refer to it as \textbf{ARMOR-FL (uncorrected)} in Section~\ref{sec:results}.5, reported exclusively as an ablation baseline against the corrected method -- not a competing method, and its numbers should be read only against the corrected version's, not Krum, FoolsGold, or FedAvg.

\section{Results}\label{sec:results}

\subsection{Aggregate accuracy across datasets}

Table~\ref{tab:master} reports mean final-round accuracy and mean last-eight-round accuracy, averaged across all twelve attack-type $\times$ malicious-fraction $\times$ non-IID-$\alpha$ combinations, for each aggregator on each dataset.

\begin{table*}[htbp]
\centering
\caption{Mean accuracy (final-round / last-8-round mean) across all twelve attack-fraction-$\alpha$ combinations per dataset}\label{tab:master}
\footnotesize
\begin{tabular}{@{}lccc@{}}
\toprule
Aggregator & CICIDS2017 & CICIDS2018 & CICIoT2023 \\
\midrule
FedAvg (undefended) & 0.727 / 0.727 & 0.825 / 0.787 & 0.413 / 0.427 \\
Krum & \textbf{0.810} / \textbf{0.781} & \textbf{0.810} / \textbf{0.809} & \textbf{0.618} / \textbf{0.623} \\
FoolsGold & 0.715 / 0.720 & 0.718 / 0.740 & 0.378 / 0.381 \\
ARMOR-FL (ours) & 0.695 / 0.714 & 0.793 / 0.731 & 0.423 / 0.422 \\
\bottomrule
\end{tabular}
\end{table*}

Krum achieves the highest aggregate mean accuracy on all three datasets, at both statistics. Against the remaining baselines the picture is mixed: ARMOR-FL trails FedAvg on CICIDS2017 (0.695/0.714 vs.\ 0.727/0.727) and CICIDS2018 (0.793/0.731 vs.\ 0.825/0.787), and is essentially tied on CICIoT2023 (0.423/0.422 vs.\ 0.413/0.427). Against FoolsGold it trails on CICIDS2017 (0.695/0.714 vs.\ 0.715/0.720), is mixed on CICIDS2018 (ahead at final round, 0.793 vs.\ 0.718, behind on last-8-mean, 0.731 vs.\ 0.740), and clearly ahead on CICIoT2023 (0.423/0.422 vs.\ 0.378/0.381). ARMOR-FL does not surpass Krum's aggregate mean on any dataset. Section~\ref{sec:discussion} discusses why Krum's winner-take-all selection appears structurally better suited to this grid than ARMOR-FL's weighted averaging, and the specific conditions under which the reverse holds.

All three datasets show meaningful oscillation between final-round and last-eight-round accuracy for every aggregator except Krum, confirming this instability is a general property of the six-of-ten random subsampling rather than specific to ARMOR-FL; Krum's winner-take-all selection is evidently more robust to this sampling noise.

\subsection{Per-attack-type breakdown}\label{sec:perattack}

Table~\ref{tab:master} aggregates across three attack types that, as Section~\ref{sec:setup} notes, differ substantially in how detectable they are from parameter-space distance alone. Table~\ref{tab:perattack} disaggregates final-round accuracy by attack type, revealing patterns the aggregate obscures.

\begin{table*}[htbp]
\centering
\caption{Final-round accuracy by dataset and attack type, mean across malicious-fraction and non-IID-$\alpha$ combinations}\label{tab:perattack}
\footnotesize
\begin{tabular}{@{}llcccc@{}}
\toprule
Dataset & Attack & FedAvg & Krum & FoolsGold & ARMOR-FL \\
\midrule
\multirow{3}{*}{CICIDS2017} & ALIE & 0.791 & \textbf{0.816} & 0.791 & 0.690 \\
 & Gaussian-noise & \textbf{0.803} & 0.757 & \textbf{0.803} & \textbf{0.803} \\
 & Label-flip & 0.586 & \textbf{0.857} & 0.549 & 0.591 \\
\midrule
\multirow{3}{*}{CICIDS2018} & ALIE & 0.824 & \textbf{0.872} & 0.831 & 0.820 \\
 & Gaussian-noise & 0.831 & \textbf{0.893} & 0.831 & 0.831 \\
 & Label-flip & \textbf{0.820} & 0.666 & 0.492 & 0.729 \\
\midrule
\multirow{3}{*}{CICIoT2023} & ALIE & \textbf{0.767} & 0.538 & 0.554 & 0.598 \\
 & Gaussian-noise & 0.023 & \textbf{0.737} & 0.023 & 0.023 \\
 & Label-flip & 0.449 & \textbf{0.579} & 0.556 & \textbf{0.647} \\
\bottomrule
\end{tabular}
\end{table*}

Three patterns stand out. First, ARMOR-FL beats Krum on label-flip on two of three datasets (CICIDS2018, CICIoT2023), consistent with the case study in Section~\ref{sec:results}.7: label-flipping is where our self-shift mechanism has genuine signal, since such an update looks central under a pure population comparison but still produces a detectable shift relative to its own pre-attack baseline. Second, under Gaussian-noise on CICIDS2017 and CICIDS2018, all four aggregators achieve nearly identical accuracy -- the injected noise is not disruptive enough at these malicious fractions to differentiate defended from undefended. Third, Gaussian-noise on CICIoT2023 produces a catastrophic collapse to $0.023$ for every aggregator except Krum, which retains $0.737$, the single largest accuracy gap anywhere in our grid, discussed as a case study in Section~\ref{sec:results}.8.

\subsection{Per-malicious-fraction and non-IID-$\alpha$ breakdown}\label{sec:perfraction}

Table~\ref{tab:perfraction} further disaggregates final-round accuracy by malicious client fraction (0.2 versus 0.4) and by non-IID Dirichlet concentration ($\alpha=0.5$ versus IID), each averaged across the remaining grid dimensions.

\begin{table*}[htbp]
\centering
\caption{Final-round accuracy by malicious fraction and by non-IID $\alpha$, mean across the remaining grid dimensions}\label{tab:perfraction}
\footnotesize
\begin{tabular}{@{}llcccc@{}}
\toprule
Dataset & Condition & FedAvg & Krum & FoolsGold & ARMOR-FL \\
\midrule
\multirow{2}{*}{CICIDS2017} & mal.\ frac.\ 0.2 & 0.823 & 0.795 & 0.791 & \textbf{0.852} \\
 & mal.\ frac.\ 0.4 & 0.630 & \textbf{0.825} & 0.638 & 0.538 \\
\midrule
\multirow{2}{*}{CICIDS2018} & mal.\ frac.\ 0.2 & 0.833 & \textbf{0.888} & 0.735 & 0.829 \\
 & mal.\ frac.\ 0.4 & \textbf{0.816} & 0.733 & 0.701 & 0.758 \\
\midrule
\multirow{2}{*}{CICIoT2023} & mal.\ frac.\ 0.2 & 0.513 & \textbf{0.728} & 0.527 & 0.552 \\
 & mal.\ frac.\ 0.4 & 0.314 & \textbf{0.508} & 0.229 & 0.294 \\
\midrule
\multirow{2}{*}{CICIDS2017} & $\alpha=0.5$ (non-IID) & 0.750 & \textbf{0.797} & 0.669 & 0.706 \\
 & IID & 0.703 & \textbf{0.823} & 0.760 & 0.683 \\
\midrule
\multirow{2}{*}{CICIDS2018} & $\alpha=0.5$ (non-IID) & 0.819 & 0.703 & 0.684 & \textbf{0.789} \\
 & IID & 0.831 & \textbf{0.917} & 0.752 & 0.798 \\
\midrule
\multirow{2}{*}{CICIoT2023} & $\alpha=0.5$ (non-IID) & 0.425 & \textbf{0.671} & 0.406 & 0.476 \\
 & IID & 0.402 & \textbf{0.566} & 0.350 & 0.370 \\
\bottomrule
\end{tabular}
\end{table*}

ARMOR-FL wins its single row on CICIDS2017 (mal.\ frac.\ 0.2) but drops sharply as malicious fraction rises to 0.4 on the same dataset (0.852 to 0.538), a larger drop than any other aggregator across the same transition -- consistent with the mechanical explanation in Section~\ref{sec:discussion}: trust-weighted averaging blends in every insufficiently-suppressed attacker, so higher malicious fraction directly increases exposure, whereas Krum's winner-take-all rule is comparatively insensitive as long as one honest client remains among the six sampled. Krum wins seven of the twelve rows outright, and ARMOR-FL's one clear win (CICIDS2018, $\alpha=0.5$) is again non-IID, reinforcing that its practical niche is non-IID rather than IID data, where population-only defenses are least reliable to begin with (Section~\ref{sec:method}).

\subsection{Detection statistics}

Table~\ref{tab:detection} reports ARMOR-FL's detection precision, recall, and benign false-exclusion rate, averaged across all twelve combinations per dataset. Precision remains low throughout (0.167 to 0.312), meaning a meaningful fraction of excluded clients are not actually malicious even after all three corrections; benign false-exclusion rates of 0.125--0.163 confirm this. Recall is lower still (0.062 to 0.229): ARMOR-FL frequently fails to exclude a truly malicious client at all, not only misfires against honest ones. This is a genuine remaining limitation, discussed further in Section~\ref{sec:limitations}.

\begin{table*}[htbp]
\centering
\caption{ARMOR-FL detection statistics, mean across all twelve combinations per dataset}\label{tab:detection}
\footnotesize
\begin{tabular}{@{}lccc@{}}
\toprule
Dataset & Precision & Recall & Benign FER \\
\midrule
CICIDS2017 & 0.312 & 0.229 & 0.125 \\
CICIDS2018 & 0.167 & 0.062 & 0.163 \\
CICIoT2023 & 0.194 & 0.125 & 0.142 \\
\bottomrule
\end{tabular}
\end{table*}

Table~\ref{tab:detectionattack} breaks this down by attack type. Gaussian-noise is detected most precisely when detected at all (up to 1.0 on CICIDS2018), but recall is frequently zero -- the aggregation-layer collapse (Section~\ref{sec:results}.2, .8) happens too fast on CICIoT2023 for the sequential test to accumulate enough evidence. ALIE, designed to evade distance-based detection \citep{baruch2019alie}, produces the highest benign false-exclusion rates on two of three datasets (0.198--0.271): an attack calibrated to blend into the honest spread also makes it harder for the population e-process to separate honest heterogeneity from the attack's camouflaged elevation.

\begin{table*}[htbp]
\centering
\caption{ARMOR-FL detection statistics by attack type, mean across malicious-fraction and non-IID-$\alpha$ combinations}\label{tab:detectionattack}
\footnotesize
\begin{tabular}{@{}llccc@{}}
\toprule
Dataset & Attack & Precision & Recall & Benign FER \\
\midrule
\multirow{3}{*}{CICIDS2017} & ALIE & 0.222 & 0.250 & 0.198 \\
 & Gaussian-noise & 0.500 & 0.250 & 0.031 \\
 & Label-flip & 0.278 & 0.188 & 0.146 \\
\midrule
\multirow{3}{*}{CICIDS2018} & ALIE & 0.000 & 0.000 & 0.271 \\
 & Gaussian-noise & 1.000 & 0.125 & 0.000 \\
 & Label-flip & 0.083 & 0.062 & 0.219 \\
\midrule
\multirow{3}{*}{CICIoT2023} & ALIE & 0.167 & 0.250 & 0.229 \\
 & Gaussian-noise & 0.000 & 0.000 & 0.062 \\
 & Label-flip & 0.333 & 0.125 & 0.135 \\
\bottomrule
\end{tabular}
\end{table*}

\subsection{Ablation and computational overhead}\label{sec:ablationresults}

\textbf{Ablation: corrected versus uncorrected ARMOR-FL.} Section~\ref{sec:ablation-protocol} describes an archived grid run under an earlier, uncorrected ARMOR-FL implementation -- before the self-shift e-process, MAD-flooring, and trust-weight corrections of Section~\ref{sec:method} -- under identical conditions to the corrected grid. Table~\ref{tab:ablation} compares the two. The corrected version reduces benign false-exclusion rate on every dataset (most sharply on CICIDS2017, 0.309 to 0.125), confirming on the full real-data grid, not only the synthetic reproductions of Section~\ref{sec:method}, that the corrections do what they were designed to do. The effect on aggregate accuracy is more mixed -- improvement on CICIoT2023 (0.395 to 0.423) and CICIDS2017 (0.689 to 0.695), a small regression on CICIDS2018 (0.805 to 0.793) -- consistent with this paper's framing: reducing false exclusions is not the same lever as maximizing accuracy, and the corrections were validated against the false-exclusion pathology specifically.

\begin{table*}[htbp]
\centering
\caption{Ablation: uncorrected versus corrected ARMOR-FL, mean across all twelve combinations per dataset}\label{tab:ablation}
\footnotesize
\begin{tabular}{@{}llccc@{}}
\toprule
Dataset & Version & Final accuracy & Precision & Benign FER \\
\midrule
\multirow{2}{*}{CICIDS2017} & Uncorrected & 0.689 & 0.288 & 0.309 \\
 & Corrected (ours) & \textbf{0.695} & \textbf{0.312} & \textbf{0.125} \\
\midrule
\multirow{2}{*}{CICIDS2018} & Uncorrected & \textbf{0.805} & \textbf{0.352} & 0.188 \\
 & Corrected (ours) & 0.793 & 0.167 & \textbf{0.163} \\
\midrule
\multirow{2}{*}{CICIoT2023} & Uncorrected & 0.395 & 0.358 & 0.208 \\
 & Corrected (ours) & \textbf{0.423} & 0.194 & \textbf{0.142} \\
\bottomrule
\end{tabular}
\end{table*}

The uncorrected version's detection precision is, counter-intuitively, higher on two of three datasets. We attribute this to a trade-off rather than a regression: the uncorrected version excludes more clients overall (hence its higher false-exclusion rate), and a subset of those additional exclusions happen to be correct, inflating precision on a smaller, more aggressively-filtered exclusion set. Section~\ref{sec:limitations} returns to this exclusion-aggressiveness-versus-accuracy tension as an unresolved open design question.

\textbf{Computational overhead.}\label{sec:overheadresults} Table~\ref{tab:overhead} reports mean wall-clock time per full run (25 rounds), averaged across all twelve combinations. Consistent with Section~\ref{sec:complexity}, Krum is the slowest of the four in aggregate, not ARMOR-FL -- its sequential-testing bookkeeping adds negligible overhead over FedAvg, while Krum's pairwise-distance computation adds a consistent, and on CICIDS2018 substantial, overhead.

\begin{table*}[htbp]
\centering
\caption{Mean wall-clock time per 25-round run (seconds), averaged across all twelve combinations per dataset}\label{tab:overhead}
\footnotesize
\begin{tabular}{@{}lcccc@{}}
\toprule
Dataset & FedAvg & Krum & FoolsGold & ARMOR-FL \\
\midrule
CICIDS2017 & 243.0 & \textbf{278.3} & 240.7 & 245.7 \\
CICIDS2018 & 285.5 & \textbf{577.5} & 358.8 & 306.1 \\
CICIoT2023 & 429.1 & \textbf{467.9} & 395.6 & 397.4 \\
\midrule
Overall mean & 319.2 & \textbf{441.2} (+38\%) & 331.7 (+4\%) & 316.4 ($-1\%$) \\
\bottomrule
\end{tabular}
\end{table*}

This is a relevant counterpoint to Section~\ref{sec:results}.1: Krum's accuracy advantage comes with a consistent computational cost -- 38\% higher aggregate wall-clock time -- that a latency-sensitive deployment may not want to pay, particularly since Section~\ref{sec:results}.3 shows the advantage narrows or reverses under the non-IID, lighter-poisoning conditions where ARMOR-FL is strongest.

Our measured numbers are for $K=10$ (six active per round), where local training dominates per-round time (Section~\ref{sec:complexity}) and the aggregation-layer cost difference is a small share of the total. We have not run experiments at larger $K$; as a theoretical projection of the aggregation step's own asymptotic cost, not an empirical extrapolation, Krum's $O(K^2d)$ cost relative to ARMOR-FL's $O(Kd)$ grows linearly with active clients per round -- tenfold at 60 active clients, hundredfold at 600 -- so at real multi-tenant SOC scale (hundreds to low thousands of clients) Krum's aggregation-layer cost disadvantage should become substantially more pronounced than the modest 38\% gap at $K=10$, though whether it matters relative to local training time at that scale is untested and left to future work.

\subsection{Case studies}\label{sec:casestudy1}

\textbf{Label-flipping under non-IID data at low malicious fraction.} The setting that most clearly illustrates when ARMOR-FL's mechanism earns its keep is label-flipping under non-IID data ($\alpha=0.5$) with 20\% malicious clients, on CICIDS2017. ARMOR-FL achieves $0.891$ final-round accuracy ($0.715$ last-8-mean), ahead of Krum ($0.811$/$0.686$), FoolsGold ($0.735$), and FedAvg ($0.799$ final-round but only $0.402$ last-8-mean). Under the same attack and fraction but IID data, the advantage disappears: Krum takes over ($0.904$/$0.836$ vs.\ ARMOR-FL's $0.864$/$0.721$). This reversal is consistent with Section~\ref{sec:method}: the shift-based gate is most effective against a genuine, attack-coincident change detected against a stable non-IID baseline, and its advantage is reported only where Table~\ref{tab:perfraction} shows it holds in aggregate too. We flag directly that this margin is a single-seed result (Section~\ref{sec:threats}): the last-8-mean margin over Krum here (0.715 vs.\ 0.686, 0.029) is smaller than FedAvg's own stability gap in this case (0.799 vs.\ 0.402, 0.397), and a different seed could narrow, erase, or reverse it. What we consider robust is the aggregate pattern that malicious fraction 0.2 under non-IID data is the one row where ARMOR-FL leads on this dataset, and the mechanistic account of why; we present this case study as illustration, not standalone evidence.

\textbf{Catastrophic Gaussian-noise collapse on CICIoT2023.}\label{sec:case2} Table~\ref{tab:perattack} reports a striking result: FedAvg, FoolsGold, and ARMOR-FL all collapse to $0.023$ under Gaussian-noise -- indistinguishable from random guessing -- while Krum retains $0.737$, the single largest gap anywhere in our grid. This is a qualitatively different failure than the one this paper otherwise addresses: not false exclusion or detection precision, but outright aggregation collapse. The per-round trace shows all three collapsing aggregators reaching near-zero accuracy within two to three post-attack rounds and never recovering. We attribute this to CICIoT2023's already-difficult multi-class task interacting with additive Gaussian perturbation at a fixed noise scale calibrated to CICIDS2017/2018 norms: on a harder task, updates plausibly occupy a narrower useful region of weight space, so a given perturbation is proportionally more destructive, and once the aggregate crosses into a sufficiently corrupted region, subsequent honest local training cannot recover within our 25-round budget for any averaging-based rule. Krum sidesteps this entirely: discarding every other contribution means a single honest client's update alone determines the aggregate, which additive noise from any other client cannot corrupt. This qualifies the cost-benefit argument of Section~\ref{sec:results}.6: Krum's overhead may be reasonable insurance against catastrophic, non-recoverable collapse. It also suggests a concrete direction for ARMOR-FL: a hard per-round trust-region bound on aggregate movement might prevent this class of collapse without Krum's winner-take-all discarding, though we have not implemented or validated this.

\section{Discussion}\label{sec:discussion}

The central empirical finding of this paper is uncomfortable in a specific way: ARMOR-FL fixes a real, previously undiagnosed failure mode exactly as intended, and Krum -- a much simpler, purely empirical rule with no statistical guarantee -- still outperforms it in aggregate on every dataset. We devote this section to unpacking why, when the advantage reverses, and what a practitioner should take from both facts.

\subsection{Why Krum wins in aggregate}\label{sec:whykrum}

Part of the explanation is mechanical. Krum discards all but one client's contribution every round, so a single successfully identified honest client keeps the aggregate clean; ARMOR-FL instead blends contributions via trust-weighted averaging, more statistically efficient when all contributors are honest but meaning one undetected malicious contributor is blended in every round it participates. Under label-flipping, where a malicious update often looks statistically central (Section~\ref{sec:method}), this cost is highest. Table~\ref{tab:perfraction} shows it growing sharply with malicious fraction: ARMOR-FL's advantage over Krum at fraction 0.2 on CICIDS2017 reverses into a deficit at 0.4, as the mechanical account predicts.

\subsection{When ARMOR-FL wins, and the CICIoT2023 difficulty confound}\label{sec:whenarmor}

The case study in Section~\ref{sec:results}.7 shows the reverse holds under the right conditions: non-IID data plus a lighter malicious fraction is where a stable per-client baseline gives the shift test real signal before a larger malicious share overwhelms the trust-weighted average. Table~\ref{tab:perattack} reinforces this by attack type: ARMOR-FL beats Krum on label-flip on two of three datasets, exactly where the self-shift mechanism is best positioned given label-flip's central-looking updates. This is the honest scope of ARMOR-FL's advantage: targeted, not universal, strongest where cumulative evidence-gathering beats an isolated heuristic -- non-IID populations, label-flipping, lighter poisoning -- before heavier malicious presence overwhelms the averaging mechanism (Section~\ref{sec:whykrum}).

CICIoT2023's uniformly lower accuracy across all four aggregators (Table~\ref{tab:master}) is worth discussing rather than treating as noise. It is a substantially harder multi-class problem than either CIC dataset -- our own centralized XGBoost baseline with full pooled-data access reached only 0.856 accuracy and 0.660 macro-F1 on CICIoT2023, against 0.999 and 0.932 on CICIDS2017 -- so much of the inter-aggregator gap likely reflects task difficulty rather than robustness specifically, and Krum's relative advantage (0.618 vs.\ ARMOR-FL's 0.423) may partly reflect a winner-take-all rule's lower sensitivity to a harder task's own noise. The catastrophic Gaussian-noise collapse (Section~\ref{sec:case2}) is consistent with this: a harder task leaves less margin for any averaging-based rule to absorb unsuppressed noise.

\subsection{Practical deployment guidance}\label{sec:deployment}

Bringing Sections~\ref{sec:results}.3, \ref{sec:results}.6, and~\ref{sec:whenarmor} together, this guidance is descriptive, not a prospective decision procedure. Attack type is adversarially chosen and not knowable in advance; non-IID-ness \emph{is} estimable in advance, e.g.\ from historical class-distribution telemetry. An operator whose threat model makes label-flipping plausible, and who values a formal, auditable false-exclusion guarantee (Section~\ref{sec:relatville}), is where ARMOR-FL earns a genuine advantage at lower cost. A deployment where malicious fraction may be high or Gaussian-noise-style corruption is plausible is better served by Krum despite its cost and absence of a guarantee -- guessing wrong carries asymmetric downside, since an averaging rule including ARMOR-FL can collapse to near-random accuracy in one scenario (Section~\ref{sec:case2}), whereas defaulting to Krum where ARMOR-FL was preferable costs only the modest gap in Table~\ref{tab:perfraction}.

\subsection{Relation to the anytime-valid inference literature}\label{sec:relatville}

A broader point: a formal Type-I guarantee (Theorem~\ref{thm:ville}) is about \emph{false-exclusion control}, not \emph{final task performance}, and these can diverge substantially. Theorem~\ref{thm:ville} is stated at $\alpha_{\mathrm{pop}}=0.05$, yet Table~\ref{tab:detection} reports realized benign false-exclusion rates of $0.125$--$0.163$, two-and-a-half to three times higher. This is not a violation -- the theorem is unconditionally true given its hypotheses -- but a consequence of those hypotheses not holding exactly for the clients the bound is meant to protect: it controls false rejection under a null of \emph{exactly} $\mu_0=1$, while a genuinely heterogeneous honest client's true mean typically sits somewhat above one, and the reported rate additionally aggregates across many clients, rounds, and grid combinations rather than one client's single-sequence rate. We state this plainly: the corrected design substantially reduces false exclusion relative to the uncorrected baseline (Section~\ref{sec:results}.5), and the theorem is mathematically correct as stated, but neither claim promises a 5\% real-world rate, and Section~\ref{sec:limitations} reports the realized rate as an open cost.

ARMOR-FL succeeds unambiguously at false-exclusion \emph{reduction} relative to its own uncorrected predecessor, but does not, on the evidence in Table~\ref{tab:master}, translate into the strongest aggregate accuracy. We think this is a useful caution for anytime-valid methods entering applied machine-learning venues, where a formal guarantee is sometimes treated as synonymous with practical superiority: it is real and valuable (auditability, principled control, no pre-committed stopping time), but a different kind of claim from task performance, and a paper introducing such a method should report both.

\section{Threats to Validity}\label{sec:threats}

We separate four categories of threat to the conclusions above, following standard empirical machine-learning reporting practice.

\textbf{Internal validity.} All results come from a single random seed per configuration; the round-to-round oscillation in Section~\ref{sec:results}.1 shows a single seed's final-round number carries non-trivial variance, and a different seed could shift some closer comparisons in Table~\ref{tab:perattack} and Table~\ref{tab:perfraction}, though the largest effects (the CICIoT2023 collapse; Krum's aggregate margin) are large enough relative to observed stability gaps to be robust to this. The case study in Section~\ref{sec:casestudy1} itself rests on a margin (0.029) smaller than the stability gap it documents for FedAvg (0.397), flagged there as needing multi-seed replication, which was outside the compute budget available (Section~\ref{sec:hyperparams}).

\textbf{External validity.} Our conclusions rest on three datasets, three attack types, two malicious-fraction levels, and two non-IID settings; extrapolating to other attack families (model-replacement, backdoor), higher fractions, other client counts, or other domains is not directly supported. The reduced sampling fractions (10\%, 2\%, 1\%, Section~\ref{sec:setup}) mean our results characterize a partial slice of each dataset, not verified against the withheld remainder.

\textbf{Construct validity.} We report final-round and last-8-round-mean accuracy, following prior convention (Section~\ref{sec:hyperparams}); this may understate performance differences on minority attack classes under imbalanced label distributions, and we did not conduct a per-class breakdown, noted as a natural extension (Section~\ref{sec:conclusion}).

\textbf{Conclusion validity.} The ablation (Section~\ref{sec:results}.5) isolates the three corrections as a bundle; we cannot attribute the false-exclusion improvement to any one in isolation, only their combination, though the synthetic reproductions in Section~\ref{sec:method} do isolate each individually. A fully factorial real-data ablation would need roughly eight times the compute budget used here and was infeasible. The regime in Section~\ref{sec:whenarmor} (non-IID, label-flipping, fraction 0.2) was also located by inspecting the grid after it completed, not predicted in advance or confirmed on a held-out slice; the mechanistic account in Section~\ref{sec:method} predicts it independently of these numbers, but a reader should weigh the finding accordingly.

\section{Limitations}\label{sec:limitations}

We state four limitations plainly, since an honest account of what ARMOR-FL does not solve is as important as the mechanism that does work.

First, ARMOR-FL cannot detect a constant attacker: a client Byzantine from round one whose bias never changes, staying below the gross-outlier threshold, produces no shift signal and is only smoothly downweighted, never hard-excluded. This is a genuine information-theoretic limit, not an implementation gap -- a constant, moderate-magnitude attacker and a heterogeneous-but-honest client are indistinguishable from behavior alone, and resolving it needs external information (a trusted validation set, cryptographic attestation) we do not assume available.

Second, label-flipping remains genuinely hard to detect via any distance-based signal, ARMOR-FL's included, since it produces an update statistically similar to an honest one -- a limitation shared with Krum, trimmed mean, and FoolsGold, not specific to our method.

Third, probation-and-reinstatement resets a client's accumulated evidence entirely, giving a reinstated attacker a fresh accumulation window before re-exclusion. We found no way to shorten this without reintroducing false exclusions for honest clients recovering from legitimate probation, and leave a more principled rule to future work.

Fourth, ARMOR-FL's own detection precision remains modest (0.167 to 0.312, Table~\ref{tab:detection}), meaning a non-trivial share of excluded clients are not actually malicious even after correction; recall is lower still (0.062 to 0.229), meaning it also frequently fails to exclude a truly malicious client. We consider this an honest residual cost rather than evidence the mechanism is not working: the corrected false-exclusion rate is dramatically lower than the uncorrected 62.5\% under IID data (Section~\ref{sec:method}), but not zero.

Fifth, the ablation (Section~\ref{sec:results}.5, Table~\ref{tab:ablation}) surfaces an unresolved tension between exclusion aggressiveness and accuracy: the uncorrected version, excluding more clients overall, achieves higher precision on two of three datasets than the corrected version. We do not have a design improving false-exclusion rate, precision, and recall simultaneously across all three datasets, and doubt one exists without external information beyond per-round behavioral distance. We view resolving this three-way trade-off as the most consequential open problem this paper leaves for future work.

\section{Conclusion}\label{sec:conclusion}

We set out to build an anytime-valid Byzantine-robust aggregator for federated intrusion detection and found that the natural approach -- comparing each client against the population via a sequential test -- carries a structural failure mode not previously diagnosed at this specificity: any anytime-valid population comparison is guaranteed to eventually exclude a persistently off-center honest client, worst precisely when the population is most homogeneous, opposite the intuitive expectation. ARMOR-FL's self-referential shift test, scale-floored population statistic, and threshold-calibrated trust weight together remove this failure mode, verified on synthetic reproductions and three real benchmarks. The method is not a uniform improvement over Krum, the strongest empirical baseline, but identifies a specific regime -- non-IID populations under lighter label-flipping poisoning -- where its formal, cumulative evidence-gathering earns a real advantage, with an anytime-valid guarantee Krum and FoolsGold do not provide. We release the full implementation, configuration grids, and reassembly-only dataset archives so this failure mode and the conditions under which our fix helps can be independently verified.

Concrete extensions follow from the limitations above. Statistically, a fully factorial ablation of the three corrections and a multi-seed repetition of the full grid (Section~\ref{sec:threats}) would sharpen the causal claims. Mechanistically, a trust-region-style bound on per-round aggregate movement (Section~\ref{sec:case2}) could plausibly close the catastrophic-collapse gap without Krum's winner-take-all discarding. Architecturally, RPCFL's cryptographic guarantee \citep{chen2026rpcfl} and ARMOR-FL's statistical guarantee address different threat surfaces (Section~\ref{sec:qualcomp}), and a combined system is unexplored. Empirically, a per-class accuracy breakdown and an extension to model-replacement and backdoor attacks (Section~\ref{sec:threats}) would test whether the regime in Section~\ref{sec:whenarmor} generalizes beyond label-flipping under non-IID data. None of these is required for the present contribution's claims to hold -- each is scoped above as a boundary of what this study supports -- but together they describe the most productive next steps.

\backmatter

\section*{Declarations}

\textbf{Funding.} No funds, grants, or other support were received for the preparation of this manuscript.

\textbf{Conflict of interest/Competing interests.} The author declares no competing interests.

\textbf{Ethics approval and consent to participate.} Not applicable. This work uses publicly available network traffic datasets generated on lab testbeds by their respective creating institutions (CICIDS2017, CICIDS2018 \citep{sharafaldin2018cicids}; CICIoT2023 \citep{neto2023ciciot2023}) rather than captured from live production networks, and involves no human subjects or personally identifiable data collected by the author.

\textbf{Dual-use considerations.} The accompanying release includes configuration files that reproduce specific poisoning attacks (label-flipping, Gaussian-noise, and the adaptive ALIE attack) against a real federated intrusion-detection backbone. We judge the marginal risk of this release to be low: all three attacks are reimplementations of already-published techniques \citep{baruch2019alie} rather than novel attack methods, the datasets are lab-testbed traffic rather than live production captures, and reproducing the attack grid requires no credentials, access, or target model beyond what the release itself provides in a fully offline simulation. We disclose this consideration explicitly rather than silently, consistent with the reporting stance taken throughout this paper.

\textbf{Consent for publication.} Not applicable.

\textbf{Data availability.} The datasets analyzed (CICIDS2017, CICIDS2018, CICIoT2023) are publicly available from the Canadian Institute for Cybersecurity and are redistributed in the accompanying code repository as checksummed archive chunks for reassembly, avoiding the need to re-download from the original hosts.

\textbf{Materials availability.} Not applicable.

\textbf{Code availability.} The full ARMOR-FL implementation, experiment configuration files, and analysis scripts are available in the accompanying code repository.

\textbf{Supplementary information.} Full per-combination result and detection-statistic tables, a reproducibility checklist, a worked numerical example, a notation summary, and baseline aggregation pseudocode are provided as Electronic Supplementary Material accompanying this article.

\textbf{Author contribution.} Q.-V.D. conceived the study, designed and implemented ARMOR-FL, ran all experiments, performed the analysis, and wrote the manuscript.

\bibliography{sn-bibliography}% common bib file

\end{document}
